\documentclass[../main.tex]{subfiles}
\begin{document}
\gdef\lecturenum{01}
\gdef\lecturetitle{Sample Lecture: Algorithms and Circuits}
\gdef\lecturedate{September 23, 2026}
\begin{center}{\Large\bfseries \coursename: \coursetitle}\\[3pt]
{\large Lecture \lecturenum --- \lecturetitle}\\[3pt]\lecturedate\end{center}
\vspace{.5em}\hrule\vspace{1em}
\section{Overview}
\begin{itemize}[leftmargin=*]\item Main topic: analyzing a simple algorithm and an RC circuit.
\item Important concept: connect equations to a practical example.
\item Connection to previous lecture: functions and basic circuit elements.\end{itemize}
\section{Core Concepts}
\begin{definitionbox}[--- Linear Function]
A linear function has the form $f(x)=ax+b$, with slope $a$ and intercept $b$.
\end{definitionbox}
For a longer derivation,
\begin{align*}y&=f(x)\\&=ax+b\\&=\boxed{ax+b}.\end{align*}
\section{Worked Example}
\begin{examplebox}[--- Matrix Determinant]
For $A=\begin{bmatrix}1&2\\3&4\end{bmatrix}$, we have
\[\det(A)=(1)(4)-(2)(3)=\boxed{-2}.\]
\end{examplebox}
\section{Algorithm / Procedure}
\begin{examplebox}[--- Sum of the First $n$ Integers]
\begin{enumerate}[leftmargin=*]\item Initialize $x\gets0$.
\item For $i=1,\ldots,n$, set $x\gets x+i$.
\item Return $x$.\end{enumerate}\end{examplebox}
\clearpage\section{Code}
\begin{intellijcode}{sum.cpp}
int sum(int n) {
    int result = 0;
    for (int i = 1; i <= n; ++i) result += i;
    return result;
}
\end{intellijcode}
\section{Circuit / Hardware}
\begin{center}\begin{circuitikz}\draw (0,0) to[V,l=$V_s$] (0,2)
to[R,l=$R$] (3,2) to[C,l=$C$] (3,0) -- (0,0);\end{circuitikz}\end{center}
\begin{importantbox}The circuit time constant is $\tau=RC$. Record the assumptions behind any model you use.\end{importantbox}
\section{Questions}
\begin{questionbox}\begin{enumerate}[leftmargin=*]
\item Why does this equation work?\item What assumptions are being made?
\item When does this method fail?\end{enumerate}\end{questionbox}
\section{Key Takeaways}
\begin{enumerate}[leftmargin=*]\item \textbf{Algorithm:} The loop adds integers from 1 through $n$.
\item \textbf{Circuit:} The product $RC$ sets the response timescale.\end{enumerate}
\section*{To Review}
\begin{itemize}[label=$\square$]\item Review: matrix determinants
\item Practice: trace the loop for $n=4$\item Ask about: circuit initial conditions\end{itemize}
\end{document}
